
\documentclass{article}
\usepackage{colortbl}
\usepackage{makecell}
\usepackage{multirow}
\usepackage{supertabular}

\begin{document}

\newcounter{utterance}

\centering \large Interaction Transcript for game `cladder', experiment `full\_v1.5\_default', episode 4 with Qwen3.5{-}9B.
\vspace{24pt}

{ \footnotesize  \setcounter{utterance}{1}
\setlength{\tabcolsep}{0pt}
\begin{supertabular}{c@{$\;$}|p{.15\linewidth}@{}p{.15\linewidth}p{.15\linewidth}p{.15\linewidth}p{.15\linewidth}p{.15\linewidth}}
    \# & \multicolumn{2}{c}{Player} && \multicolumn{2}{c}{Game Master} \\
    \hline

    \theutterance \stepcounter{utterance}  
    & & & \multicolumn{4}{p{0.6\linewidth}}{
        \cellcolor[rgb]{0.9,0.9,0.9}{
            \makecell[{{p{\linewidth}}}]{
                \texttt{\tiny{[P1$\langle$GM]}}
                \texttt{You are an expert in causal inference. The following question is not a typical commonsense query, but rather a meticulously designed question created by a professor specializing in causal inference, intended to assess the students' mastery of the course content.} \\
\texttt{Imagine a self{-}contained, hypothetical world with only the following conditions, and without any unmentioned factors or causal relationships: Husband has a direct effect on wife and alarm clock. Wife has a direct effect on alarm clock. For husbands that don't set the alarm and wives that don't set the alarm, the probability of ringing alarm is 11\%. For husbands that don't set the alarm and wives that set the alarm, the probability of ringing alarm is 60\%. For husbands that set the alarm and wives that don't set the alarm, the probability of ringing alarm is 46\%. For husbands that set the alarm and wives that set the alarm, the probability of ringing alarm is 92\%. For husbands that don't set the alarm, the probability of alarm set by wife is 61\%. For husbands that set the alarm, the probability of alarm set by wife is 1\%. Does husband positively affect alarm clock through wife?} \\
\texttt{Start your answer with 'Yes' or 'No', followed by additional reasoning or evidence to support your explanation.} \\
            }
        }
    }
     \\ \\

    \theutterance \stepcounter{utterance}  
    & \multicolumn{4}{p{0.6\linewidth}}{
        \cellcolor[rgb]{0.561,0.737,0.561}{
            \makecell[{{p{\linewidth}}}]{
                \texttt{\tiny{[P1$\rangle$GM]}}
                \texttt{No} \\
\\ 
\texttt{**Reasoning:**} \\
\\ 
\texttt{To determine if the husband positively affects the alarm clock *through* the wife, we must calculate the **Indirect Effect** of the husband on the alarm clock via the wife. In causal inference, this is defined as the difference in the outcome (alarm ringing) when the husband changes his action, while the wife's action is held fixed at the level it would have been under the husband's original action (or simply comparing the counterfactuals where the wife's action is mediated).} \\
\\ 
\texttt{Let us define the variables:} \\
\texttt{*   \$H\$: Husband sets alarm (\$H=1\$) vs. does not (\$H=0\$).} \\
\texttt{*   \$W\$: Wife sets alarm (\$W=1\$) vs. does not (\$W=0\$).} \\
\texttt{*   \$Y\$: Alarm rings (\$Y=1\$) vs. does not (\$Y=0\$).} \\
\\ 
\texttt{The question asks for the effect of \$H\$ on \$Y\$ *through* \$W\$. This corresponds to the path \$H \textbackslash{}to W \textbackslash{}to Y\$. Mathematically, the indirect effect is often calculated as:} \\
\texttt{\$\$ \textbackslash{}text\{Indirect Effect\} = E{[}Y | H=1, W=1{]} {-} E{[}Y | H=0, W=1{]} \$\$} \\
\texttt{*(Note: This specific formulation isolates the change in the outcome due to the change in the mediator \$W\$ caused by \$H\$, holding the mediator at} \\
            }
        }
    }
    & & \\ \\

    \theutterance \stepcounter{utterance}  
    & & & \multicolumn{2}{p{0.3\linewidth}}{
        \cellcolor[rgb]{0.9,0.9,0.9}{
            \makecell[{{p{\linewidth}}}]{
                \texttt{\tiny{[GM$|$GM]}}
                \texttt{No} \\
            }
        }
    }
    & & \\ \\

    \theutterance \stepcounter{utterance}  
    & & & \multicolumn{2}{p{0.3\linewidth}}{
        \cellcolor[rgb]{0.9,0.9,0.9}{
            \makecell[{{p{\linewidth}}}]{
                \texttt{\tiny{[GM$|$GM]}}
                \texttt{no} \\
            }
        }
    }
    & & \\ \\

    \theutterance \stepcounter{utterance}  
    & & & \multicolumn{2}{p{0.3\linewidth}}{
        \cellcolor[rgb]{0.9,0.9,0.9}{
            \makecell[{{p{\linewidth}}}]{
                \texttt{\tiny{[GM$|$GM]}}
                \texttt{game\_result = WIN} \\
            }
        }
    }
    & & \\ \\

\end{supertabular}
}

\end{document}
